\section{Regular immersions and normal flatness}
\label{section:IV.19}

\oldpage[IV]{185}
The present paragraph is devoted, on the one hand, to a study of \emph{regular immersions} (\sref{IV.16.9.2}) $Y\to X$ of preschemes, notably in the case where $X$ and $Y$ are flat over a same prescheme $S$; on the other hand, it gives complements on M-regular suites and normal flatness, generalising notably results on flatness established by H.~Hironaka in his theory of resolution of singularities \cite{IV-35}.

\subsection{Properties of regular immersions.}
\label{subsection:IV.19.1}

\begin{proposition}[19.1.1]
\label{IV.19.1.1}
Let $X$ be a locally Noetherian prescheme, $j:Y\to X$ an immersion, $y$ a regular point of $Y$.
For $j$ to be a regular immersion at the point $y$, it is necessary and sufficient that $X$ be regular at the point $j(y)$.
\end{proposition}

Taking into account (\sref{IV.16.9.10}), one is reduced to proving the

\begin{corollary}[19.1.2]
\label{IV.19.1.2}
Let $A$ be a local Noetherian ring, $\mathfrak{J}$ an ideal of $A$; suppose that the quotient ring $B=A/\mathfrak{J}$ is regular.
For $A$ to be regular, it is necessary and sufficient that the ideal $\mathfrak{J}$ be regular.
\end{corollary}

Indeed, if $A$ is regular, we know ($\mathbf{0}$, \sref{IV.17.1.9}) that $\mathfrak{J}$ is generated by a part of a system of parameters of $A$, hence $\mathfrak{J}$ is regular.
Inversely, if $\mathfrak{J}$ is regular, and if $(x_i)_{1\leq i\leq r}$ is an A-regular suite generating $\mathfrak{J}$, $(x_i)$ forms part of a system of parameters of $A$ ($\mathbf{0}$, \sref{IV.16.4.1}), hence one deduces from ($\mathbf{0}$, \sref{IV.17.1.7}) that $A$ is regular.

\begin{definition}[19.1.3]
\label{IV.19.1.3}
Let $X$ be a prescheme, $Y$ a sub-prescheme of $X$, $U$ an open of $X$ containing $Y$ and such that $Y$ is defined by an Ideal $\mathscr{J}$ of $\sh{O}_U$; suppose that $\mathscr{J}/\mathscr{J}^2$ is an $\sh{O}_U/\mathscr{J}$-Module locally free (which is the case if $Y$ is quasi-regularly immersed in $X$).
For every $y\in Y$, we call the \emph{transverse codimension} of $Y$ in $X$ at the point $y$ and write $\codim^*_y(Y,X)$ the rank of the $(\sh{O}_{y}/\mathfrak{J}_y)$-module libre $\mathscr{J}_y/\mathscr{J}_y^2$.
If $f:Z\to X$ is an immersion with image $Y$, we similarly call \emph{transverse codimension} of $f$ at the point $z\in Z$ the transverse codimension in $X$ at the point $f(z)$ of the sub-prescheme $Y$.
\end{definition}

\begin{proposition}[19.1.4]
\label{IV.19.1.4}
Let $X$ be a locally Noetherian prescheme, $Y$ a sub-prescheme of $X$ regularly immersed; then, for every $y\in Y$, one has
\begin{equation}
\label{IV.19.1.4.1}
  \codim^*_y(Y,X) = \codim_y(Y,X).
\end{equation}
\end{proposition}

By virtue of (\sref{IV.5.1.3.2}), $\codim_y(Y,X)$ is equal to the least value of the dimensions of the local rings $A=\sh{O}_{X,z}$, where $z$ is a generic point of an irreducible component of $Y$ containing $y$; as such a point $z$ is contained in every neighbourhood of $y$ in $X$, $\mathscr{J}_z/\mathscr{J}_z^2$ is an A-module free of rank $n=\codim^*_y(Y,X)$; it amounts to seeing that $\dim(A)=n$.
But, since $z$ is a maximal point of the sub-prescheme $Y$ defined by $\mathscr{J}$, one has $\dim(\sh{O}_{X,z}/\mathscr{J}_z)=0$, hence $\mathscr{J}_z$ is an \emph{ideal of definition} of the local Noetherian ring $A$; moreover, the hypothesis that $\mathscr{J}$ is quasi-regular implies that $\mathscr{J}_z^m/\mathscr{J}_z^{m+1}$ is an $(A/\mathscr{J}_z)$-module free of rank $\binom{m+n-1}{n-1}$; if $r$ is the length of the Artinian ring $A/\mathscr{J}_z$, the length of $\mathscr{J}_z^m/\mathscr{J}_z^{m+1}$ is thus $r\binom{m+n-1}{n-1}$; the Hilbert-Samuel polynomial of $A$ for the $\mathscr{J}_z$-adic filtration is thus of degree $n$, which proves the proposition ($\mathbf{0}$, \sref{IV.16.2.3}).

By virtue of (\sref{IV.19.1.4}), we shall henceforth say ``codimension'' instead of ``transverse codimension'' when it concerns a sub-prescheme $Y$ regularly immersed in $X$ (\emph{even} when $X$ is not locally Noetherian).

\begin{proposition}[19.1.5]
\label{IV.19.1.5}
\oldpage[IV]{186}
\emph{(i)} For an immersion $j:Y\to X$ to be open, it is necessary and sufficient that it be regular and of codimension $0$ at every point.

\emph{(ii)} Let $f:Y\to X$ be a regular immersion (resp.\ quasi-regular), $g:X'\to X$ a flat morphism, $Y'=Y\times_X X'$, $f'=f_{(X')}:Y'\to X'$.
Then $f'$ is a regular immersion (resp.\ quasi-regular), and for every $y'\in Y'$, if $y$ is the projection of $y'$ in $Y$, the codimension of $f'$ at the point $y'$ is equal to the codimension of $f$ at the point $y$.
Inversely, if $f$ is an immersion and $g:X'\to X$ a faithfully flat and quasi-compact morphism, then, if $f'$ is a quasi-regular immersion, so is $f$.

\emph{(iii)} If $f:Y\to X$ and $g:Z\to Y$ are regular immersions, then $f\circ g:Z\to X$ is a regular immersion, and for every $z\in Z$, the codimension of $f\circ g$ at the point $z$ is the sum of the codimension of $g$ at the point $z$ and of the codimension of $f$ at the point $g(z)$.
Moreover, the suite of canonical homomorphisms (\sref{IV.16.2.7.1})
\begin{equation}
\label{IV.19.1.5.1}
  0\to g^*(\sh{N}_{Y/X})\to\sh{N}_{Z/X}\to\sh{N}_{Z/Y}\to 0
\end{equation}
is exact, and for every $z\in Z$, there exists a neighbourhood of $z$ in $Z$ in which the restriction homomorphisms of this suite form a split exact suite.

\emph{(iv)} Let $X$ be a locally Noetherian prescheme, $f:Y\to X$ and $g:Z\to Y$ two immersions, $z$ a point of $Z$.
The following conditions are equivalent:
\begin{enumerate}[label=\emph{\alph*)}]
  \item $g$ is a regular immersion at the point $z$ and $f$ a regular immersion at the point $g(z)$.
  \item $g$ and $f\circ g$ are regular immersions at the point $z$.
  \item $f\circ g$ is a regular immersion at the point $z$, $f$ is a regular immersion at the point $g(z)$, and the suite of canonical homomorphisms (\sref{IV.19.1.5.1}), restricted to a sufficiently small open neighbourhood of $z$ in $Z$, is exact and split.
\end{enumerate}

\emph{(v)} Let $X$ be a locally Noetherian prescheme, $f:Y\to X$, $g:Z\to Y$ two immersions, $z$ a point of $Z$; suppose that $y=g(z)$ is a regular point of $Y$ and $x=f(g(z))$ a regular point of $X$ (which implies, by (\sref{IV.19.1.1}), that $f$ is a regular immersion at the point $g(z)$); \emph{then}, for $f\circ g$ to be a regular immersion at the point $z$, it is necessary and sufficient that $g$ be a regular immersion at the point $z$.
\end{proposition}

\begin{corollary}[19.1.6]
\label{IV.19.1.6}
\oldpage[IV]{188}
Let $A$ be a local Noetherian ring, $\mathfrak{J}$, $\mathfrak{K}$ two ideals of $A$ contained in the maximal ideal $\mathfrak{m}$ and such that $\mathfrak{J}\subset\mathfrak{K}$, $A'=A/\mathfrak{J}$, $\mathfrak{K}'=\mathfrak{K}/\mathfrak{J}$.
Then:
\begin{enumerate}[label=\emph{(\roman*)}]
  \item If $\mathfrak{J}$ and $\mathfrak{K}'$ are regular ideals, so is $\mathfrak{K}$.
  \item If $\mathfrak{K}$ and $\mathfrak{K}'$ are regular ideals, so is $\mathfrak{J}$.
  \item Suppose the rings $A$ and $A'=A/\mathfrak{J}$ are regular.
  For the ideal $\mathfrak{K}$ to be regular, it is necessary and sufficient that $\mathfrak{K}'$ be so.
\end{enumerate}
\end{corollary}

Assertion (i) is a particular case of (\sref{IV.19.1.5}[iii]).
To prove (ii), consider the surjective canonical homomorphism $u:\mathfrak{K}/\mathfrak{K}^2\to\mathfrak{K}'/\mathfrak{K}'^2=\mathfrak{K}/(\mathfrak{K}^2+\mathfrak{J})$.
Since by hypothesis $\mathfrak{K}/\mathfrak{K}^2$ and $\mathfrak{K}'/\mathfrak{K}'^2$ are $(A/\mathfrak{K})$-modules free, whose ranks we denote $p$ and $q$ respectively, the kernel of $u$ is an $(A/\mathfrak{K})$-module \emph{projective} (Bourbaki, \emph{Alg.~comm.}, chap.~II, \S~2, n\textsuperscript{o}~2, prop.~4), hence \emph{free} of rank $p$ since $A/\mathfrak{K}$ is a local Noetherian ring ($\mathbf{0_{III}}$, \sref{IV.10.1.3}); in other words, there exists a base of $\mathfrak{K}/\mathfrak{K}^2$ whose first $p$ elements form a base of $(\mathfrak{K}^2+\mathfrak{J})/\mathfrak{K}^2$.
This means (by virtue of Nakayama's lemma) that there exists a system of generators $(f_i)_{1\leq i\leq p+q}$ of $\mathfrak{K}$, forming a regular suite in $A$, such that the suite $(f_i)_{1\leq i\leq p}$ is formed of elements of $\mathfrak{J}$; it remains to show that this last suite generates $\mathfrak{J}$.
For this, denoting by $\mathfrak{J}'\subset\mathfrak{J}$ the ideal it generates, one sees that one is reduced to proving the following lemma:

\begin{lemma}[19.1.6.1]
\label{IV.19.1.6.1}
Let $B$ be a local Noetherian ring with maximal ideal $\mathfrak{n}$, $(g_j)_{1\leq j\leq q}$ a regular suite in $B$ of elements of $\mathfrak{n}$, $\mathfrak{b}$ the ideal that it generates, $\mathfrak{c}\subset\mathfrak{b}$ an ideal such that the canonical images of the $g_j$ in $B/\mathfrak{c}$ form a regular suite in this ring.
Then $\mathfrak{c}=0$.
\end{lemma}

The lemma being trivial for $q=0$, one argues by induction on $q$.
The induction hypothesis, applied to the ring $B/g_1 B$ and to the canonical images of the $g_j$ ($2\leq j\leq q$), shows that $\mathfrak{c}\subset g_1\mathfrak{b}$.
Let $\mathfrak{b}$ be the ideal of $B$ of elements $x\in B$ such that $g_1 x\in\mathfrak{c}$; one therefore has $\mathfrak{c}=g_1\mathfrak{b}$, but since by hypothesis $g_1$ is regular in $B/\mathfrak{c}$, one has necessarily $\mathfrak{b}=\mathfrak{c}$, hence $\mathfrak{c}=g_1\mathfrak{c}$.
Since $\mathfrak{c}$ is of finite type and $g_1$ is contained in the radical of $B$, Nakayama's lemma shows that $\mathfrak{c}=0$.

Finally, taking into account (\sref{IV.16.9.10}), it is immediate that (v), as well as the equivalence of \emph{a)} and \emph{b)} in (iv), result from

\begin{remark}[19.1.7]
\label{IV.19.1.7}
\emph{(i)} We ignore whether the composition of two quasi-regular immersions is quasi-regular.

\emph{(ii)} In an annelated space in local rings $X$, for every Ideal $\mathscr{J}$ of $\sh{O}_X$, $\sh{O}_X/\mathscr{J}$ is of finite type and thus has closed support in $X$ ($\mathbf{0_I}$, \sref{IV.5.2.2}).
One can thus define as in (\sref[I]{I.4.1.3} and \sref[I]{I.4.2.1}) the notions of closed subspace (locally closed subspace) of $X$ defined by an Ideal $\mathscr{J}$ of a sheaf $\sh{O}_U$, where $U$ is an open of $X$, and the notion of immersion.
The definitions of quasi-regular and regular immersions, of transverse codimension, and the results of (\sref{IV.19.1.4}[i] to [iv]) remain valid without modification, assuming in the second assertion of (i) and in (iv) that the sheaf $\sh{O}_X$ is coherent and the local rings $\sh{O}_{X,x}$ are Noetherian; as for (ii), one must \emph{define} $Y'$ as the closed subspace defined by the $\sh{O}_X$-Ideal $g^*(\mathscr{J})$; the proofs are unchanged.
\end{remark}

\begin{env}[19.1.8]
\label{IV.19.1.8}
\oldpage[IV]{189}
We have already pointed out (notably in (\sref{IV.16.9.6}[ii])) that in non-Noetherian rings, the notion of ``regular suite'' of elements of a ring does not possess the properties that would render it usable.
Recent researches seem to show that the notion which should in this case be substituted for that of regular suite is that of a suite $\mathbf{f}=(f_1,\ldots,f_n)$ having the property that the Koszul complex $K_\bullet(\mathbf{f},M)$ (\sref[III]{III.1.1.3}) has \emph{zero homology in degrees $>0$} (cf.\ (\sref[III]{III.1.1.4})); see in particular A.~Grothendieck, S\'eminaire de g\'eom\'etrie alg\'ebrique, 1966, to appear shortly.
\end{env}

\subsection{Transversally regular immersions.}
\label{subsection:IV.19.2}

\oldpage[IV]{190}

\begin{definition}[19.2.1]
\label{IV.19.2.1}
Let $u:X\to S$ be a morphism of preschemes, $\mathscr{F}$ an $\sh{O}_X$-Module quasi-coherent.
One says that a suite $(f_i)_{1\leq i\leq n}$ of sections of $\mathscr{F}$ above $X$ is \emph{transversally $\mathscr{F}$-regular relatively to $S$} (or relatively to $u$) if the suite $(f_i)$ is $\mathscr{F}$-regular and such that the $\sh{O}_X$-Modules $\mathscr{F}/(\sum_{i<j}f_i\mathscr{F})$ are S-flat for $0\leq i\leq n$.
One says that an Ideal quasi-coherent $\mathscr{J}$ of $\sh{O}_X$ is \emph{transversally regular relatively to $S$} (or relatively to $u$) at a point $x\in\operatorname{Supp}(\sh{O}_X/\mathscr{J})$ if there exists an open neighbourhood $U$ of $x$ and a finite suite $(f_i)$ of sections of $\mathscr{J}$ which is $\sh{O}_U$-regular relatively to $S$ and which generates $\mathscr{J}|U$.

In an A-algebra $B$, one says that an ideal $\mathfrak{J}$ is \emph{transversally regular relatively to $A$} if $\mathscr{J}=\widetilde{\mathfrak{J}}$ is transversally regular relatively to $S=\operatorname{Spec}(A)$ at every point of $V(\mathfrak{J})$ in $\operatorname{Spec}(B)$.
\end{definition}

\begin{definition}[19.2.2]
\label{IV.19.2.2}
Let $S$ be a prescheme, $X$, $Y$ two S-preschemes, $u:Y\to X$ an S-morphism which is an immersion.
One says that $u$ is an immersion \emph{transversally regular relatively to $S$} at a point $y\in Y$ if there exists a neighbourhood $V$ of $u(y)$ in $X$ such that the sub-prescheme of $X$ associated to $u$ is defined by an Ideal of $\sh{O}_U$ transversally regular relatively to $S$ at the point $u(y)$.

It is clear that the set of points of $Y$ where an immersion is transversally regular relatively to $S$ is \emph{open} in $Y$; if it is equal to $Y$, one says simply that $u$ is an immersion transversally regular relatively to $S$.

If $S\to T$ is a flat morphism, then, if $\mathscr{J}$ is an Ideal of $\sh{O}_X$ transversally regular relatively to $S$ at $x\in X$, it is so also relatively to $T$ ($\mathbf{0_I}$, \sref{IV.6.2.1}).
A $Y$-immersion $Y\to X$ transversally regular relatively to $S$ at a point $y\in Y$ relatively to $S$ is therefore also a T-immersion transversally regular at this point.
\end{definition}

\begin{proposition}[19.2.3]
\label{IV.19.2.3}
Let $S$ be a prescheme, $X$, $Y$ two S-preschemes, $u:Y\to X$ an S-morphism which is an immersion transversally regular relatively to $S$ at a point $y\in Y$.
Then, for every change of base $S'\to S$, if one sets $X'=X\times_S S'$, $Y'=Y\times_S S'$, the immersion $u'=u_{(S')}:Y'\to X'$ is transversally regular relatively to $S'$ at every point $y'\in Y'$ above $y$.
\end{proposition}

This follows immediately from the definition and ($\mathbf{0}$, \sref{IV.15.1.15}).

\begin{proposition}[19.2.4]
\label{IV.19.2.4}
Let $S$ be a prescheme, $X$, $Y$ two S-preschemes, $u:Y\to X$ an S-morphism which is an immersion.
Suppose, either that $S$ and $X$ are locally Noetherian, or that the structural morphisms $g:X\to S$, $h:Y\to S$ are locally of finite presentation.
Let $y$ be a point of $Y$, $s=g(x)=h(y)$.
The following conditions are equivalent:
\begin{enumerate}[label=\emph{\alph*)}]
  \item $u$ is transversally regular relatively to $S$ at the point $y$.
  \item The morphisms $g$ and $h$ are flat at the neighbourhoods of the points $u(y)$ and $y$ respectively, and if one sets $X_s=g^{-1}(s)$, $Y_s=h^{-1}(s)$, the immersion $u_s:Y_s\to X_s$ is regular at the point $y$.
  \item[b')] The morphisms $g$ and $h$ are flat at the neighbourhoods of the points $u(y)$ and $y$ respectively, and, for every change of base $S'\to S$, if one sets $X'=X\times_S S'$, $Y'=Y\times_S S'$, the immersion $u'=u_{(S')}:Y'\to X'$ is regular at every point of $Y'$ above $y$.
  \item The morphism $h$ is flat at the neighbourhood of the point $y$, and the immersion $u$ is regular at the neighbourhood of the point $y$.
  \item[c')] The morphism $h$ is flat at the neighbourhood of the point $y$, and the immersion $u$ is quasi-regular at the neighbourhood of the point $y$.
\end{enumerate}
\end{proposition}

\begin{remark}[19.2.5]
\label{IV.19.2.5}
\oldpage[IV]{191}
\emph{(i)} When in condition $b)$ one suppresses the hypothesis that $h$ is flat at the neighbourhood of $y$, the condition thus modified no longer implies $a)$.
It suffices to take for $S$ the spectrum of a local ring which is not a field, for $s$ its closed point, and $Y=X_s$.

\emph{(ii)} Suppose that $Y$ is a closed sub-prescheme of $X$ defined by an Ideal quasi-coherent $\mathscr{J}$ of $\sh{O}_X$.
One has proved in the course of the proof of (\sref{IV.19.2.4}) that, when the equivalent conditions of this proposition are satisfied, then, for every system $(f_i)$ of sections of $\mathscr{J}$ from a neighbourhood of $y$ in $X$, such that the images of $f_i$ in $\mathscr{J}_y/(\mathfrak{m}_s\sh{O}_{X,y}+\mathscr{J}_y^2)$ form a base of this $(\sh{O}_{X,y}/(\mathfrak{m}_s\sh{O}_{X,y}+\mathscr{J}_y))$-module (which is free by virtue of the hypothesis $b)$), there exists an open neighbourhood $U$ of $y$ in $X$ such that the $f_i|U$ satisfy the conditions of the definition (\sref{IV.19.2.1}).
\end{remark}

\begin{corollary}[19.2.6]
\label{IV.19.2.6}
\oldpage[IV]{192}
Suppose, either that $S$ and $X$ are locally Noetherian, or that $g$ and $h$ are locally of finite presentation.
Suppose furthermore that the fibres $X_s$ and $Y_s$ are regular preschemes at the points $u(y)$ and $y$ respectively, and that the morphisms $g$ and $h$ are flat at the neighbourhoods of the points $u(y)$ and $y$ respectively.
Then the immersion $Y\to X$ is transversally regular at the point $x$.
\end{corollary}

Indeed, it follows from (\sref{IV.19.1.1}) that the immersion $Y_s\to X_s$ is regular at the point $y$, and it suffices to apply (\sref{IV.19.2.4}[$b)$]).

\begin{proposition}[19.2.7]
\label{IV.19.2.7}
\emph{(i)} Every open S-immersion is transversally regular relatively to $S$.

\emph{(ii)} Let $f:Y\to X$ be an S-immersion, $g:S'\to S$ a morphism, and set $X'=X_{(S')}$, $f'=f_{(S')}:Y'\to X'$.
If $f$ is an immersion transversally regular relatively to $S$, $f'$ is transversally regular relatively to $S'$.
The converse is true if furthermore $g$ is faithfully flat, and if moreover $X$ and $Y$ are locally of finite presentation over $S$.

\emph{(iii)} If $f:Y\to X$ and $g:Z\to Y$ are immersions transversally regular relatively to $S$, so is $f\circ g:Z\to X$.

\emph{(iv)} Let $X$ be an S-prescheme, $f:Y\to X$, $g:Z\to Y$ two S-immersions.
Suppose, either that $S$ and $X$ are locally Noetherian, or that $X$, $Y$ and $Z$ are locally of finite presentation over $S$.
Then, if $g$ and $f\circ g$ are transversally regular relatively to $S$ at the point $z\in Z$, $f$ is transversally regular relatively to $S$ at the point $g(z)$.

\emph{(v)} Under the general conditions of (iv), suppose that $X$ and $Y$ are S-flat, $Y_s$ regular at the point $g(z)$ and $X_s$ regular at the point $f(g(z))$; \emph{then}, if $f\circ g$ is transversally regular at the point $z$, so is $g$.
\end{proposition}

\begin{proposition}[19.2.8]
\label{IV.19.2.8}
\oldpage[IV]{193}
Let $X$ be an S-prescheme, $Y$ a closed sub-prescheme of $X$ defined by an Ideal quasi-coherent $\mathscr{J}$ of $\sh{O}_X$, $u:Y\to X$ the canonical injection.
If $u$ is transversally regular relatively to $S$, the $\sh{O}_X$-Modules $\mathscr{J}^n$ and $\mathscr{J}^n/\mathscr{J}^m$ ($0\leq n<m$) are S-flat at all points of $Y$.
For every change of base $S'\to S$, if $X'=X_{(S')}$, $Y'=Y_{(S')}$ and $u'=u_{(S')}$, one has $\mathscr{G}_n(u')=\mathscr{G}_n(u)\otimes_{\sh{O}_S}\sh{O}_{S'}$ as isomorphism for every $n>0$.
\end{proposition}

Indeed, $\sh{O}_X$ and $\sh{O}_X/\mathscr{J}$ are by hypothesis S-flat at points of $Y$, and $\mathscr{J}^n/\mathscr{J}^{n+1}$ is thus S-flat at points of $Y$.
The first assertion follows then from ($\mathbf{0_I}$, \sref{IV.6.1.2}) and the second from (\sref{IV.11.2.9.2}).

The following proposition generalises and makes precise (\sref{IV.17.16.1}):

\begin{proposition}[19.2.9]
\label{IV.19.2.9}
Let $f:X\to S$ be a morphism, $x$ a point of $X$, $s=f(x)$.
Suppose $f$ is of finite presentation at the point $x$, and is a Cohen-Macaulay morphism at the point $x$ (\sref{IV.6.8.1}).
Then there exists a sub-prescheme $Y$ of $X$ containing $x$ and having the following properties:
\begin{enumerate}[label=\emph{(\roman*)}]
  \item The canonical injection $j:Y\to X$ is transversally regular relatively to $S$.
  \item The morphism $g=f\circ j:Y\to S$ is de Cohen-Macaulay (which is equivalent to saying, taking into account (i), that all the fibres $g^{-1}(g(y))$ of $g$ are Cohen-Macaulay preschemes).
  \item The point $x$ is a maximal point of $Y_s=g^{-1}(s)$.
\end{enumerate}

Set $X_s=f^{-1}(s)$; the ring $\sh{O}_{X_s,x}$ is Cohen-Macaulay; let $(t_i)_{1\leq i\leq n}$ be a system of parameters of this ring, which is therefore a regular suite ($\mathbf{0}$, \sref{IV.16.5.7}).
There is a neighbourhood open $U$ of $x$ in $X$ and sections $h_i$ ($1\leq i\leq n$) of $\sh{O}_X$ above $U$ whose canonical images $(h_i)_x$ in the quotient ring $\sh{O}_{X_s,x}$ coincide with $t_i$; and $\sh{O}_{Y_s,x}$ is flat over $\sh{O}_{S,s}$; by virtue of (\sref{IV.11.3.1}) $f$ and $g$ are flat in a neighbourhood of $x$ and it follows from (\sref{IV.19.2.4}) that $j$ is an immersion transversally regular relatively to $S$ (replacing $U$ by a suitable open neighbourhood of $x$ in $X$).
Finally, the ring $\sh{O}_{Y_s,x}$, quotient of $\sh{O}_{X_s,x}$ by the ideal generated by the $t_i$, is Artinian by construction ($\mathbf{0}$, \sref{IV.16.3.6}), so $x$ is a maximal point of $Y_s$.
Finally, the ring $\sh{O}_{Y_s,x}$, being Artinian, is a Cohen-Macaulay ring ($\mathbf{0}$, \sref{IV.16.5.1}), and, replacing $Y$ by its intersection with a suitable open neighbourhood $V$ of $x$ in $X$, one can suppose that $g$ is a Cohen-Macaulay morphism, by virtue of (\sref{IV.12.1.7}[iii]).

In particular, one recovers the existence of ``quasi-sections'' \emph{flat} of a morphism $f:X\to S$ at a point $x\in X_s$ which is \emph{closed} in $X_s$ and such that $f$ is flat at the point $x$ and that $\sh{O}_{X_s,x}$ is a Cohen-Macaulay ring (\sref{IV.17.16.1}); one sees moreover that the immersion $Y\to X$ is \emph{transversally regular relatively to $S$}.
This last complement is also valid for the quasi-section \'etale defined in (\sref{IV.17.16.3}[i]).
\end{proposition}

\subsection{Relative complete intersections (flat case).}
\label{subsection:IV.19.3}

\oldpage[IV]{194}

\begin{definition}[19.3.1]
\label{IV.19.3.1}
Let $A$ be a local Noetherian ring.
One says that $A$ is a ring of \emph{complete intersection} \emph{(or equivalently, a ring of absolute complete intersection, if there is risk of confusion)} if the completion $\hat{A}$ is isomorphic to the quotient of a local Noetherian regular complete ring $B$ by a regular ideal (i.e.\ (\sref{IV.16.9.7}) an ideal generated by a regular suite of elements of $B$).

Every regular local ring is a ring of complete intersection.

One says that a locally Noetherian prescheme $X$ is a \emph{complete intersection (absolute)} at a point $x\in X$ if the local ring $\sh{O}_{X,x}$ is a ring of complete intersection.
\end{definition}

\begin{proposition}[19.3.2]
\label{IV.19.3.2}
Let $B$ be a local Noetherian regular ring, $\mathfrak{J}$ an ideal of $B$.
For $A=B/\mathfrak{J}$ to be a ring of complete intersection, it is necessary and sufficient that the ideal $\mathfrak{J}$ be regular (i.e.\ (\sref{IV.16.9.7}) generated by a regular suite of elements of $B$).
\end{proposition}

One can write $\hat{A}=\hat{B}/\mathfrak{J}\hat{B}$ and since $\hat{B}$ is a B-module faithfully flat and $B$ Noetherian, $\mathfrak{J}\hat{B}$ is regular if and only if $\mathfrak{J}$ is (\sref{IV.19.1.5}[ii]).
This shows (by virtue of ($\mathbf{0}$, \sref{IV.17.1.5}[ii])) that the statement is sufficient, and on the other hand that to prove it is necessary, one can restrict to the case where $A$ and $B$ are \emph{complete}.
By hypothesis, there exists then a local Noetherian complete \emph{regular} ring $B'$ such that $A$ is isomorphic to a quotient $B'/\mathfrak{J}'$, where $\mathfrak{J}'$ is a regular ideal of $B'$.
Consider then the fibred product $B'=B\times_A B''$ for the surjective homomorphisms $B\to A$, $B'\to A$ ($\mathbf{0}$, \sref{IV.18.1.2}); we shall first prove the following lemma:

\begin{lemma}[19.3.2.1]
\label{IV.19.3.2.1}
Let $A$, $E$, $F$ be three rings, $f:E\to A$, $g:F\to A$ two homomorphisms, and set $G=E\times_A F$ ($\mathbf{0}$, \sref{IV.18.1.2}).
\begin{enumerate}[label=\emph{(\roman*)}]
  \item If $f$ is surjective and if $E$ and $F$ are Noetherian rings, so is $G$.
  \item If $A$, $E$ and $F$ are local rings and if $f$ and $g$ are local homomorphisms, $G$ is local, and the canonical homomorphisms $G\to E$, $G\to F$ are local.
  \item If the conditions of (i) and (ii) are satisfied, if $g$ is surjective, and if $E$ and $F$ are complete local Noetherian rings, so is $G$.
\end{enumerate}
\end{lemma}

This lemma being established, it follows from the Cohen theorem ($\mathbf{0}$, \sref{IV.19.8.8}) that $B''$ is isomorphic to a quotient of a local Noetherian complete \emph{regular} ring $C$.
It follows then from (\sref{IV.19.1.2}) that the immersion $\operatorname{Spec}(B')\to\operatorname{Spec}(C)$ is regular; since by hypothesis, we conclude from (\sref{IV.19.1.5}[iii]) that the immersion $\operatorname{Spec}(A)\to\operatorname{Spec}(C)$ also writes as the composed $\operatorname{Spec}(A)\to\operatorname{Spec}(B)\to\operatorname{Spec}(C)$.
But since $B$ is regular by hypothesis, we conclude from (\sref{IV.19.1.5}[iii]) that this immersion is also composed as $\operatorname{Spec}(A)\to\operatorname{Spec}(B)\to\operatorname{Spec}(C)$; hence it is also the same of $\operatorname{Spec}(A)\to\operatorname{Spec}(B)$ by (\sref{IV.19.1.2}); hence it is also of the same of $\operatorname{Spec}(A)\to\operatorname{Spec}(B)$ by (\sref{IV.19.1.5}[iv]).

\begin{corollary}[19.3.3]
\label{IV.19.3.3}
Let $X$ be a locally Noetherian prescheme locally immersible in a regular prescheme.
Then the set of $x\in X$ where $X$ is a complete intersection is open in $X$ and contains the set of regular points of $X$.
\end{corollary}

One can indeed restrict to the case where $X=\operatorname{Spec}(B/\mathfrak{J})$, where $B$ is a regular ring and $\mathfrak{J}$ an ideal of $B$.
The set of $\mathfrak{p}\supset\mathfrak{J}$ in $\operatorname{Spec}(B)$ such that $\mathfrak{J}_\mathfrak{p}$ is regular in $B_\mathfrak{p}$ is open (\sref{IV.16.9.5}), and this is the set of points of $\operatorname{Spec}(B)$ where $X$ is a complete intersection by virtue of (\sref{IV.19.3.2}).

\begin{corollary}[19.3.4]
\label{IV.19.3.4}
Let $k$ be a field, $k'$ an extension of $k$, $X'=X\otimes_k k'$.
Let $x'$ be a point of $X'$, $x$ its projection in $X$.
For $X$ to be a complete intersection at $x$, it is necessary and sufficient that $X'$ be a complete intersection at $x'$.
\end{corollary}

The question being local on $X$, one can suppose that $X=\operatorname{Spec}(A)$, where $A$ is a $k$-algebra of finite type, hence a quotient of a polynomial algebra $k[T_1,\ldots,T_n]$, so that $X$ is a closed sub-prescheme of the regular scheme $Y=\operatorname{Spec}(k[T_1,\ldots,T_n])$ ($\mathbf{0}$, \sref{IV.17.3.7}).
If one sets $Y'=Y\otimes_k k'=\operatorname{Spec}(k'[T_1,\ldots,T_n])$, $Y'$ is also a regular scheme.
Since $\operatorname{Spec}(k')$ is faithfully flat over $\operatorname{Spec}(k)$, it follows from (\sref{IV.19.1.5}[ii]) and the fact that these are locally Noetherian preschemes that the immersion $Y\to X$ is regular at a point $x$ if and only if the immersion $Y'\to X'$ is regular at every point $x'$ above $x$.
The conclusion follows from the criterion (\sref{IV.19.3.2}).

\begin{corollary}[19.3.5]
\label{IV.19.3.5}
\oldpage[IV]{195}
Let $A$, $C$ be two local Noetherian rings, $\rho:A\to C$ a local homomorphism, $B=C/\mathfrak{J}$ a quotient of $C$, $k$ the residue field of $A$.
Suppose that $\rho$ makes $C$ a flat A-module, and that the ring $C\otimes_A k$ is regular.
Then the following conditions are equivalent:
\begin{enumerate}[label=\emph{\alph*)}]
  \item The ideal $\mathfrak{J}$ is transversally regular relatively to $A$.
  \item $B$ is a flat A-module, and $B\otimes_A k$ is a ring of complete intersection.
\end{enumerate}
\end{corollary}

By virtue of (\sref{IV.19.2.4}), the condition $a)$ is equivalent to saying that $B$ is a flat A-module and that $\mathfrak{J}\otimes_A k$ is a regular ideal of $C\otimes_A k$ (by virtue of the flatness of $B$ over $A$ ($\mathbf{0_I}$, \sref{IV.6.1.2})); as the ring $C\otimes_A k$ is regular, it follows at once, by virtue of (\sref{IV.19.3.2}), that $\mathfrak{J}\otimes_A k$ is a regular ideal of $C\otimes_A k$, or that $B\otimes_A k=(C\otimes_A k)/(\mathfrak{J}\otimes_A k)$ is a ring of complete intersection; whence the corollary.

\begin{definition}[19.3.6]
\label{IV.19.3.6}
\oldpage[IV]{196}
Let $f:X\to S$ be a flat morphism locally of finite presentation.
One says that $X$ is a \emph{complete intersection relatively to $S$} at a point $x$, if the fibre $f^{-1}(f(x))$ is a complete intersection (absolute) at the point $x$.
One says that $X$ is a \emph{complete intersection relatively to $S$}, or that the morphism $f$ is a \emph{morphism of complete intersection}, if $X$ is a complete intersection relatively to $S$ at each of its points.
\end{definition}

\begin{proposition}[19.3.7]
\label{IV.19.3.7}
Let $g:X\to S$, $h:Y\to S$ be two flat morphisms locally of finite presentation, $f:Y\to X$ an S-immersion, $y$ a point of $Y$, $x=f(y)$, $s=g(x)=h(y)$.
Suppose that the fibre $X_s=g^{-1}(s)$ is regular at the point $x$.
Then, for $f$ to be an immersion transversally regular at the point $x$, it is necessary and sufficient that $Y$ be a complete intersection relatively to $S$ at the point $y$.
\end{proposition}

To say that $f$ is transversally regular at the point $y$ amounts, by virtue of (\sref{IV.19.2.4}), to saying that $f_s:Y_s\to X_s$ is a regular immersion at the point $y$.
But since $X_s$ is regular at the point $x=f(y)$, this is equivalent, by virtue of (\sref{IV.19.3.2}), to saying that $Y_s$ is a complete intersection at the point $y$, whence the proposition.

\begin{corollary}[19.3.8]
\label{IV.19.3.8}
Let $f:X\to S$ be a flat morphism locally of finite presentation.
The set $U$ of $x\in X$ at which $X$ is a complete intersection relatively to $S$ is open in $X$.
If moreover $f$ is proper, the set $E$ of $s\in S$ such that $X_s=f^{-1}(s)$ is a complete intersection is an open subset of $S$.
\end{corollary}

Since $E=Y-f(X-U)$, the second assertion follows from the first and the fact that $f$ is a closed map.
The first assertion is local on $X$, hence one can suppose that $X$ is a closed sub-prescheme of a prescheme of the form $Z=S[T_1,\ldots,T_n]$ ($T_i$ indeterminates).
The latter being smooth over $S$ (\sref{IV.17.3.8}), its fibres $Z_s$ are regular, and it amounts, by virtue of (\sref{IV.19.3.7}), to saying that $x\in U$ if and only if the immersion $X\to Z$ is transversally regular at the point $x$; one concludes thus from the corollary by (\sref{IV.19.2.2}).

\begin{proposition}[19.3.9]
\label{IV.19.3.9}
\emph{(i)} Every open immersion is a morphism of complete intersection.

\emph{(ii)} Let $f:X\to S$ be a flat morphism locally of finite presentation which is a morphism of complete intersection.
For every change of base $g:S'\to S$, $f'=f_{(S')}:X_{(S')}\to S'$ is a morphism of complete intersection.
The converse is true if $g$ is faithfully flat and quasi-compact.

\emph{(iii)} If $f:X\to Y$ and $g:Y\to Z$ are two flat morphisms locally of finite presentation which are morphisms of complete intersection, so is $g\circ f:X\to Z$.
\end{proposition}

Assertion (i) follows immediately from the definition (\sref{IV.19.3.1}).
To prove (ii), we first note that if $f$ is flat and locally of finite presentation, so is $f'$ (\sref{IV.2.1.4}), the converse being true if $g$ is faithfully flat ((\sref{IV.2.5.1}) and (\sref{IV.2.7.1})); moreover for every $s'\in S'$, if $s=g(s')$, it is equivalent to say that $f'^{-1}(s')$ is a complete intersection or that $f'^{-1}(s')$ is a complete intersection; whence (ii), taking into account of (\sref{IV.19.3.4}).
Finally, under the hypotheses of (iii), $g\circ f$ is flat and locally of finite presentation.
By definition (\sref{IV.19.3.6}), one is thus reduced to the case where $Z$ is the spectrum of a field, and to prove that the local rings $\sh{O}_{X,x}$ are rings of complete intersection, which is contained in the following more general result:

\begin{corollary}[19.3.10]
\label{IV.19.3.10}
\oldpage[IV]{197}
Let $f:X\to S$ be a flat morphism locally of finite presentation, $x$ a point of $X$, $s=f(x)$.
If $\sh{O}_{S,s}$ is a local Noetherian ring of complete intersection, so is $\sh{O}_{X,x}$.
\end{corollary}

\begin{proof}
One can evidently restrict to the case where $S=\operatorname{Spec}(A)$ with $A=\sh{O}_{S,s}$.
We show moreover that one can restrict to the case where the local ring $A$ is \emph{complete}.
Indeed, set $A'=\hat{A}$, $X'=X\otimes_A A'$, and let $s'$ be the unique closed point of $S'=\operatorname{Spec}(A')$, which is the unique point of $S'$ above $s$; if $f'=f_{(S')}:X'\to S'$, the fibre $f'^{-1}(s')$ is canonically isomorphic to $f^{-1}(s)$ (\sref[I]{I.3.6.4}); there is thus a single point $x'\in X'$ above $x$ and $s'$, and one has by suite $\sh{O}_{X',x'}=\sh{O}_{X,x}\otimes_A A'$.
One deduces that the local rings $\sh{O}_{X,x}$ and $\sh{O}_{X',x'}$ have \emph{isomorphic completions}, for, if $E$ is a local ring, the separated completion $(E\otimes_A\hat{A})^\wedge$ of the ring $E\otimes_A\hat{A}$ equipped with the topology induced by the $\mathfrak{m}$-preadic topology of $G$, is isomorphic to $\hat{E}\otimes_A\hat{A}\cong\hat{E}$, hence to the separated completion $\hat{E}$ of $E$ ($\mathbf{0_I}$, \sref{IV.7.7.1}).
By virtue of the definition of complete intersection (\sref{IV.19.3.1}), it follows that $\sh{O}_{X,x}$ is a ring of complete intersection at the same time as $\sh{O}_{X',x'}$.

Suppose therefore $A$ complete; one can moreover suppose that $X=\operatorname{Spec}(B)$, $B$ being a quotient of a polynomial algebra $C=A[T_1,\ldots,T_r]$.
Since polynomial rings over fields are regular rings ($\mathbf{0}$, \sref{IV.17.3.7}), it follows from (\sref{IV.19.3.7}) that the immersion $X\to Z=\operatorname{Spec}(C)$ can be supposed transversally regular relatively to $S$, so that one can suppose that $B=C/\mathfrak{J}$, where $\mathfrak{J}$ is generated by a regular suite $(g_i)_{1\leq i\leq n}$ of elements of $C$.
On the other hand, since $A$ is complete, one can by hypothesis write $A=A'/\mathfrak{K}$, where $A'$ is a local regular ring and $\mathfrak{K}$ an ideal of $A'$ ($\mathbf{0}$, \sref{IV.19.8.8}), and the hypothesis that $A$ is a ring of complete intersection implies that $\mathfrak{K}$ is generated by a regular suite $(f_j)_{1\leq j\leq m}$ (\sref{IV.19.3.2}).
In the polynomial ring $C'=A'[T_1,\ldots,T_r]$, the elements $f_j$ ($1\leq j\leq m$) form a regular suite $\mathfrak{K}'=\mathfrak{K}[T_1,\ldots,T_r]$ ($\mathbf{0}$, \sref{IV.15.1.4}); and the suite formed of the $f_j$ ($1\leq j\leq m$) and the $g_i$ ($1\leq i\leq n$) is regular in $C'$; if $\mathfrak{J}'$ is the ideal of $C'$ that it generates, one has $B=C'/\mathfrak{J}'$, whence the conclusion by virtue of (\sref{IV.19.3.2}), since $C'$ is a regular ring ($\mathbf{0}$, \sref{IV.17.3.7}).
\end{proof}

\subsection{Application: criteria of regularity and smoothness for blow-ups.}
\label{subsection:IV.19.4}

\oldpage[IV]{198}

\begin{env}[19.4.1]
\label{IV.19.4.1}
Let $X$ be a prescheme, $Y$ a closed sub-prescheme of $X$, defined by an Ideal quasi-coherent $\mathscr{J}$ of $\sh{O}_X$.
Recall (\sref[II]{II.8.1.3}) that the X-scheme $X'$ obtained by blowing up $Y$ is the prescheme
\[
  X'=\operatorname{Proj}(\mathscr{S}),\quad\text{where}\quad\mathscr{S}=\bigoplus_{n\geq 0}\mathscr{J}^n.
\]
If $\mathscr{J}$ is of finite type, the structural morphism $f:X'\to X$ is projective.
Without further hypothesis on $\mathscr{J}$, the closed sub-prescheme $Y'=f^{-1}(Y)$ of $X'$ is defined by the Ideal quasi-coherent $\mathscr{J}\sh{O}_{X'}$ of $\sh{O}_{X'}$, canonically isomorphic to $\sh{O}_{X'}(1)$; more precisely, one has an exact suite
\begin{equation}
\label{IV.19.4.1.1}
  0\to\sh{O}_{X'}(1)\xrightarrow{\widetilde{u}_0}\sh{O}_{X'}\to\sh{O}_{Y'}\to 0
\end{equation}
where $\widetilde{u}_0$ is the image of $\widetilde{u}_0$ (\sref[II]{II.8.1.7} and \sref[II]{II.8.1.8}).
Since, in a neighbourhood of a point of $X'$, the $\sh{O}_{X'}$-Module $\sh{O}_{X'}(1)$ is free of rank $1$, $\mathscr{J}\sh{O}_{X'}$ is generated in such a neighbourhood by an invertible section of $\sh{O}_{X'}$ and $\mathscr{J}\sh{O}_{X'}/\mathscr{J}^2\sh{O}_{X'}$ is therefore an $\sh{O}_{X'}/\mathscr{J}\sh{O}_{X'}$-Module free of rank $1$.
This proves the first assertion of the following lemma:
\end{env}

\begin{lemma}[19.4.2]
\label{IV.19.4.2}
The canonical immersion $Y'\to X'$ is regular and of codimension $1$, and $Y'$ is canonically Y-isomorphic to the Y-scheme $\operatorname{Proj}(\operatorname{gr}^*_{\mathscr{J}}(\sh{O}_X))$.
\end{lemma}

The second assertion follows from (\sref[II]{II.3.5.3}) applied to the canonical injection $Y\to X$, taking into account that by definition, $\sh{O}_Y$ is isomorphic to $\sh{O}_X/\mathscr{J}$ and $\mathscr{J}^n/\mathscr{J}^{n+1}$.

\begin{corollary}[19.4.3]
\label{IV.19.4.3}
If the canonical immersion $Y\to X$ is quasi-regular, the restriction $g:Y'\to Y$ of the morphism $f$ is smooth.
\end{corollary}

Indeed, $\sh{N}_{Y/X}=\mathscr{J}/\mathscr{J}^2$ is then a $\sh{O}_Y$-Module locally free of finite rank, hence $Y'$ is Y-isomorphic to $\mathbf{P}(\sh{N}_{Y/X})$ (\sref{IV.16.9.8}), hence $Y'$ is smooth over $Y$ (\sref{IV.17.3.9}).

We shall see later (chap.~V) that $g$ is also smooth when the immersion $Y\to X$ ``presents only ordinary singularities''.

\begin{proposition}[19.4.4]
\label{IV.19.4.4}
With the notations of (\sref{IV.19.4.1}), suppose that $X$ is locally Noetherian and that the morphism $g:Y'\to Y$, restriction of $f$, is smooth.
Let $x'$ be a point of $Y'$, $x$ its image in $Y$.
For $Y$ to be regular at the point $x$, it is necessary and sufficient that $Y'$ be regular at the point $x'$; when this is the case, $X'$ is also regular at the point $x'$.
\end{proposition}

The first assertion follows from (\sref{IV.17.5.8}) and the second from (\sref{IV.19.1.1}) and (\sref{IV.19.4.2}).

\begin{corollary}[19.4.5]
\label{IV.19.4.5}
Under the general hypotheses of (\sref{IV.19.4.4}), suppose that $Y$ is regular at the point $x$.
Then, for every generalisation $x_1$ of $x$ in $X$ not belonging to $Y$, $X$ is regular at the point $x_1$.
If $\operatorname{Reg}(X)$ is open ((\sref{IV.6.12}), for example if $X$ is excellent (\sref{IV.7.8.6}[iii]))
\end{corollary}

\begin{proposition}[19.4.6]
\label{IV.19.4.6}
\oldpage[IV]{199}
Let $g:X\to S$ be a morphism, $Y$ a closed sub-prescheme of $X$ defined by an Ideal quasi-coherent $\mathscr{J}$, $X'$ the X-scheme obtained by blowing up $Y$, $Y'$ the inverse image of $Y$ in $X'$.
The following conditions are equivalent:
\begin{enumerate}[label=\emph{\alph*)}]
  \item $\operatorname{gr}^*_{\mathscr{J}}(\sh{O}_X)$ is a $g$-flat $\sh{O}_X$-Algebra.
  \item For every $n\geq 0$, the $\sh{O}_X$-Module $\sh{O}_X/\mathscr{J}^{n+1}$ is $g$-flat (in other words, the $n$-th infinitesimal neighbourhood $Y^{(n)}$ of $Y$ in $X$ (\sref{IV.16.1.2}) is S-flat).
  \item For every change of base $S_1\to S$ and every integer $n\geq 1$, if one sets $X_1=X\times_S S_1$, the canonical homomorphism $\mathscr{J}^n\otimes_{\sh{O}_S}\sh{O}_{S_1}\to\mathscr{J}^n\sh{O}_{X_1}$ is bijective.
\end{enumerate}
When this is the case, $Y'$ is S-flat, and for every change of base $S_1\to S$, if $Y_1$ is the inverse image of $Y$ in $X_1$, the prescheme $X'_1$ obtained by blowing up $Y_1$ in $X_1$ is canonically isomorphic to $X'\times_S S_1$.
\end{proposition}

\begin{corollary}[19.4.7]
\label{IV.19.4.7}
\oldpage[IV]{200}
Let $g:X\to S$ be a morphism locally of finite presentation, $Y$ a closed sub-prescheme of $X$ such that the composed morphism $Y\to X\to S$ is locally of finite presentation; suppose moreover that $Y$ is S-flat and that $\sh{O}_X$ is normally flat along $Y$ (\sref{IV.11.3.4}).
Then (with the notations of (\sref{IV.19.4.1})), the morphism $X'\to X$ is of finite presentation, the canonical immersion $Y'\to X'$ is transversally regular of codimension $1$ relatively to $S$, and the set of points of $X$ (resp.\ $X'$) where $X$ is S-flat (resp.\ $Y'$) is a neighbourhood of $Y$ (resp.\ $Y'$).
\end{corollary}

\begin{proposition}[19.4.8]
\label{IV.19.4.8}
Suppose verified the hypotheses of (\sref{IV.19.4.7}) and suppose in addition that the morphism $Y'\to Y$ is smooth.
Let $x'$ be a point of $Y'$, $x$ its image in $Y$.
For $Y$ to be smooth over $S$ at the point $x$, it is necessary and sufficient that $Y'$ be smooth over $S$ at the point $x'$; when this is the case, $X'$ is also smooth over $S$ at the point $x'$.
\end{proposition}

\begin{corollary}[19.4.9]
\label{IV.19.4.9}
\oldpage[IV]{201}
Under the general hypotheses of (\sref{IV.19.4.8}), suppose in addition that $Y$ is smooth over $S$ at the point $x$.
Then, for every generalisation $x_1$ of $x$ in $X$ not belonging to $Y$, $X$ is smooth over $S$ at the point $x_1$.
If $Y$ is smooth over $S$, there exists an open neighbourhood $U$ of $Y$ in $X$ such that $X$ is smooth over $S$ at the points of $U-Y$.
\end{corollary}

\begin{remark}[19.4.10]
\label{IV.19.4.10}
In (\sref{IV.19.4.6}[ii]), replacing the hypothesis that $X$ and $Y$ are locally of finite presentation over $S$ by that of $S$ and $X$ (hence also $Y$) being \emph{locally Noetherian}.
Then it is clear that $\bigoplus_{n\geq 0}\mathscr{J}^n$ is an $\sh{O}_X$-Algebra of finite type, that the morphism $X'\to X$ is of finite type, and it follows still from (\sref{IV.19.2.4}) applied to locally Noetherian preschemes, that the immersion $Y'\to X'$ is transversally regular relatively to $S$ and that $X'$ is S-flat at points of $Y'$.
\end{remark}

\begin{proposition}[19.4.11]
\label{IV.19.4.11}
Let $X$ be a prescheme, $\mathscr{E}$ an $\sh{O}_X$-Module quasi-coherent, $u:\mathscr{E}\to\sh{O}_X$ a homomorphism of $\sh{O}_X$-Modules, $\mathscr{J}=u(\mathscr{E})$ the Ideal quasi-coherent of $\sh{O}_X$ image of $u$, $Y$ the closed sub-prescheme of $X$ defined by $\mathscr{J}$.
Let $X'$ be the X-scheme obtained by blowing up $Y$.
On the other hand, let $P=\mathbf{P}(\mathscr{E})$ be the projective fibre defined by $\mathscr{E}$ (\sref[II]{II.4.1.1}), $p:P\to X$ the structural morphism, $\alpha_1^{\#}:p^*(\mathscr{E})\to\sh{O}_P(1)$ the canonical homomorphism (\sref[II]{II.4.1.5.1}) and $\mathscr{H}$ its kernel, so that one has the exact suite
\[
  0\to\mathscr{H}\to p^*(\mathscr{E})\xrightarrow{\alpha_1^{\#}}\sh{O}_P(1)\to 0.
\]
Finally, let $\mathscr{K}$ the Ideal quasi-coherent of $\sh{O}_P$ image of the restriction $v:\mathscr{H}\to\sh{O}_P$ of $p^*(u)$, and let $Z$ be the closed sub-prescheme of $P$ defined by $\mathscr{K}$.
Then:
\begin{enumerate}[label=\emph{(\roman*)}]
  \item The image of $X'$ by the closed immersion $j:X'\to\mathbf{P}(\mathscr{E})$ corresponding to the surjective graduated homomorphism $\mathbf{S}_{\sh{O}_X}(\mathscr{E})\to\mathbf{S}_{\sh{O}_X}(\mathscr{J})\to\bigoplus_{n\geq 0}\mathscr{J}^n$ (\sref[II]{II.3.6.2}) is a closed sub-prescheme of $Z$.
  \item If $\mathscr{E}$ is an $\sh{O}_X$-Module locally free of rank $m$, $\mathscr{H}$ is an $\sh{O}_P$-Module locally free of rank $m-1$.
  \item Suppose that $X$ is locally Noetherian, that $\mathscr{E}$ is locally free of finite type, that the sub-prescheme $Y$ is regular, and that at every point $x\in Y$, the canonical immersion $i:Y\to X$ is regular and of codimension equal to $\operatorname{rg}_{\sh{O}_x}(\mathscr{E}_x)$ (which we express by saying that the homomorphism $u$ is regular at the point $x$).
  Then the image of $X'$ by the closed immersion $j:X'\to P$ is the closed sub-prescheme $Z$ of $P$; at all points of $X'-Y'$ ($Y'$ the inverse image of $Y$ in $X'$) the immersion $j$ is transversally regular relatively to $X$; moreover at all points $x'$ of $Y'$, $X'$ is regular, the immersion $j$ is regular and of codimension $\operatorname{rg}_{\sh{O}_{P,z}}(\mathscr{K}_z)$ where $z=j(x')$ (in other words the homomorphism $v$ is also regular at the point $z$).
\end{enumerate}
\end{proposition}

\subsection{Criteria of M-regularity.}
\label{subsection:IV.19.5}

\oldpage[IV]{204}
We revisit here, in complementing them, the criteria for a suite of elements of a ring $A$ to be M-regular or M-quasi-regular ($M$ being an A-module), already studied in ($\mathbf{0}$, \sref{IV.15.1}).

\begin{theorem}[19.5.1]
\label{IV.19.5.1}
Let $A$ be a ring, $\mathbf{f}=(f_i)_{1\leq i\leq n}$ a suite of elements of $A$, $M$ an A-module.
Consider the following properties:
\begin{enumerate}[label=\emph{\alph*)}]
  \item The suite $\mathbf{f}$ is M-regular.
  \item $\mathrm{H}_i(\mathbf{f},M)=0$ for every $i>0$ (\sref[III]{III.1.1.3}).
  \item[b')] $\mathrm{H}_1(\mathbf{f},M)=0$.
  \item The suite $\mathbf{f}$ is M-quasi-regular.
\end{enumerate}
Then one has the implications
\[
  a)\Rightarrow b)\Rightarrow b')\quad\text{and}\quad a)\Rightarrow c).
\]

Set $\mathfrak{J}=\sum_{i=1}^n f_i A$.
If the modules $M_i=M/(\sum_{j=1}^i f_j M)$ are separated for the $\mathfrak{J}$-preadic topology (resp.\ if every quotient module of a sub-module of $M$ is separated for the $\mathfrak{J}$-preadic topology), then one also has the implication $c)\Rightarrow a)$ (resp.\ $b')\Rightarrow a)$).
\end{theorem}

\begin{corollary}[19.5.2]
\label{IV.19.5.2}
\oldpage[IV]{205}
Suppose that $A$ is a Noetherian ring, that $M$ is an A-module of finite type, and that the $f_i$ ($1\leq i\leq n$) belong to the radical of $A$.
Then the four conditions \emph{a)}, \emph{b)}, \emph{b')}, \emph{c)} of (\sref{IV.19.5.1}) are equivalent.
\end{corollary}

Indeed, one knows in this case ($\mathbf{0_I}$, \sref{IV.7.3.5}) that every A-module of finite type is separated for the $\mathfrak{J}$-preadic topology.

\begin{proposition}[19.5.3]
\label{IV.19.5.3}
Let
\[
  0\to M'\to M\to M''\to 0
\]
be an exact suite of A-modules, $\mathbf{f}=(f_i)_{1\leq i\leq n}$ a suite M-regular, $\mathfrak{J}=\sum_{i=1}^n f_i A$.
Consider the following properties:
\begin{enumerate}[label=\emph{\alph*)}]
  \item The suite $\mathbf{f}$ is $M''$-regular.
  \item The $\mathfrak{J}$-preadic filtration of $M'$ is induced by the $\mathfrak{J}$-preadic filtration of $M$ (in other words, the canonical homomorphism $\operatorname{gr}^*_{\mathfrak{J}}(M')\to\operatorname{gr}^*_{\mathfrak{J}}(M)$ is injective).
  \item The canonical homomorphism $M'/(\sum_{i=1}^n f_i M')\to M/(\sum_{i=1}^n f_i M)$ is injective.
\end{enumerate}
Then one has the implications $a)\Rightarrow b)\Rightarrow c)$, and $a)$ implies furthermore that the suite $\mathbf{f}$ is $M'$-regular.
If every quotient of a sub-module of $M''$ is separated for the $\mathfrak{J}$-preadic topology, the conditions \emph{a)}, \emph{b)} and \emph{c)} are equivalent.
\end{proposition}

\begin{env}[19.5.4]
\label{IV.19.5.4}
\oldpage[IV]{206}
In the remainder of this section, we keep the notations $A$, $\mathbf{f}$, $M$, $\mathfrak{J}$ and suppose moreover $M$ equipped with a decreasing filtration $(M_k)$ formed of sub-A-modules, with $M_0=M$.
Recall ($\mathbf{0}$, \sref{IV.15.1.5}) that one defines then on $M$ a second decreasing filtration formed of the sub-modules
\[
  M'_k=M_k+\mathfrak{J}M_{k-1}+\ldots+\mathfrak{J}^{k-1}M_1+\mathfrak{J}^k M_0
\]
and that, if $\operatorname{gr}_.(M)$ and $\operatorname{gr}'_.(M)$ are the graded A-modules associated to the filtrations $(M_k)$ and $(M'_k)$ respectively, one defines a surjective graded homomorphism of degree $0$
\[
  \psi_M:(\operatorname{gr}_.(M)\otimes_A(A/\mathfrak{J}))[\mathrm{T}_1,\ldots,\mathrm{T}_n]\to\operatorname{gr}'_.(M)
\]
($\mathbf{0}$, \sref{IV.15.1.5.2}).
Recall also that one has the surjective canonical homomorphism ($\mathbf{0}$, \sref{IV.15.1.1.1})
\[
  \varphi_{\operatorname{gr}_.(M)}:(\operatorname{gr}_.(M)\otimes_A(A/\mathfrak{J}))[\mathrm{T}_1,\ldots,\mathrm{T}_n]\to\operatorname{gr}^*_{\mathfrak{J}}(\operatorname{gr}_.(M)).
\]
\end{env}

\begin{theorem}[19.5.5]
\label{IV.19.5.5}
With the notations of (\sref{IV.19.5.4}), consider the following properties:
\begin{enumerate}[label=\emph{\alph*)}]
  \item The suite $\mathbf{f}$ is $\operatorname{gr}_.(M)$-regular.
  \item The canonical homomorphism $\psi_M$ is bijective.
  \item The canonical homomorphism $\varphi_{\operatorname{gr}_.(M)}$ is bijective (in other words ($\mathbf{0}$, \sref{IV.15.1.7}), the suite $\mathbf{f}$ is $\operatorname{gr}_.(M)$-quasi-regular).
\end{enumerate}
One has the implications
\[
  a)\Rightarrow b)\Rightarrow c).
\]
Furthermore, if, for every integer $k\geq 0$, the quotient modules of sub-modules of $\operatorname{gr}_k(M)$ are separated for the $\mathfrak{J}$-preadic topology (which is the case when $A$ is Noetherian, the $f_i$ are in the radical of $A$ and the $\operatorname{gr}_k(M)$ are A-modules of finite type), \emph{then the conditions \emph{a)}, \emph{b)} and \emph{c)} are equivalent}.
\end{theorem}

\begin{remark}[19.5.6]
\label{IV.19.5.6}
\oldpage[IV]{209}
\emph{(i)} The results of this section remain valid when, instead of an A-module $M$ and elements $f_i\in A$, one takes an object $M$ of an arbitrary abelian category, and $n$ endomorphisms $f_i$ of $M$ which commute two by two; the proofs adapt without difficulty to this more general case.
The hypothesis of separation for an A-module $N$ relatively to the $\mathfrak{J}$-preadic topology should here be replaced by the hypothesis that every sub-object of $N$ contained in all the $\sum f_i^r(N)$ ($r$ an arbitrary integer) should be zero.
The same remark applies to the results of \sref{IV.19.7}.

\emph{(ii)} Suppose that $A$ is a graded ring with degrees bounded below $\geq 0$.
If $M$ (resp.\ each $\operatorname{gr}_p(M)$) is a graded module with degrees bounded below and if the $f_i$ contain no homogeneous component of degree $0$, the hypothesis of separation of ($\mathbf{0}$, \sref{IV.15.1.9}) is \emph{ipso facto} verified for $M$ (resp.\ each $\operatorname{gr}_p(M)$), and one sees thus that in this case one has the implication $c)\Rightarrow a)$ in (\sref{IV.19.5.1}) (resp.\ (\sref{IV.19.5.5})).

\emph{(iii)} Recall that, without hypothesis of separation, the implication $c)\Rightarrow b)$ is no longer valid even for $n=1$ ($\mathbf{0}$, \sref{IV.15.1.12}[iii]).
\end{remark}

\subsection{Regular suites relatively to a filtered quotient module.}
\label{subsection:IV.19.6}

\oldpage[IV]{209}

\begin{env}[19.6.1]
\label{IV.19.6.1}
Let $A$ be a ring, $\mathbf{f}=(f_1,\ldots,f_n)$ a finite suite of elements of $A$, $\mathfrak{J}=f_1 A+\ldots+f_n A$ the ideal that it generates, and consider an exact suite of A-modules
\[
  0\to R\xrightarrow{j} N\xrightarrow{p} M\to 0
\]
where we suppose $N$ equipped with a decreasing filtration $(N_k)$ formed of sub-A-modules, with $N_0=N$.
One sets $R_k=R\cap N_k$ (filtration induced by $(N_k)$), $M_k=p(N_k)$ (quotient filtration of $(N_k)$) and one designates by $\operatorname{gr}_.(N)$, $\operatorname{gr}_.(R)$ and $\operatorname{gr}_.(M)$ the graded A-modules associated to these three filtrations.
One sets on the other hand for every $k$,
\[
  N'_k=N_k+\mathfrak{J}N_{k-1}+\ldots+\mathfrak{J}^{k-1}N_1+\mathfrak{J}^k N_0
\]
so that $M'_k=p(N'_k)=M_k+\mathfrak{J}M_{k-1}+\ldots+\mathfrak{J}^{k-1}M_1+\mathfrak{J}^k M_0$; one sets on the other hand $R'_k=R\cap N'_k$ (filtration \emph{induced} by $(N'_k)$), and one notes $\operatorname{gr}'_.(N)$, $\operatorname{gr}'_.(M)$ and $\operatorname{gr}'_.(R)$ the graded A-modules associated to these three filtrations; one has thus a commutative diagram of exact suites
\begin{equation}
\label{IV.19.6.1.1}
\begin{gathered}
\xymatrix{
  0\ar[r]&\operatorname{gr}_.(R)\ar[r]^-{\operatorname{gr}(j)}\ar[d]&\operatorname{gr}_.(N)\ar[r]^-{\operatorname{gr}(p)}\ar[d]&\operatorname{gr}_.(M)\ar[r]\ar[d]&0\\
  0\ar[r]&\operatorname{gr}'_.(R)\ar[r]^-{\operatorname{gr}'(j)}&\operatorname{gr}'_.(N)\ar[r]^-{\operatorname{gr}'(p)}&\operatorname{gr}'_.(M)\ar[r]&0
}
\end{gathered}
\end{equation}
where the vertical arrows are the canonical homomorphisms of a graded module associated to a filtration to the graded module associated to a finer filtration.
\end{env}

\begin{env}[19.6.2]
\label{IV.19.6.2}
\oldpage[IV]{210}
One has defined in ($\mathbf{0}$, \sref{IV.15.1.5.2}) surjective graded homomorphisms of degree $0$
\begin{align*}
  \psi_N &: (\operatorname{gr}_.(N)\otimes_A(A/\mathfrak{J}))[\mathrm{T}_1,\ldots,\mathrm{T}_n]\to\operatorname{gr}'_.(N)\\
  \psi_M &: (\operatorname{gr}_.(M)\otimes_A(A/\mathfrak{J}))[\mathrm{T}_1,\ldots,\mathrm{T}_n]\to\operatorname{gr}'_.(M)\\
  \psi_R &: (\operatorname{gr}_.(R)\otimes_A(A/\mathfrak{J}))[\mathrm{T}_1,\ldots,\mathrm{T}_n]\to\operatorname{gr}''_.(R)
\end{align*}
where the first two are deduced from the two last vertical arrows of (\sref{IV.19.6.1.1}).
We denote by $\psi'_R$ the homomorphism composed $\gamma\circ\psi_R$.
\end{env}

\begin{proposition}[19.6.3]
\label{IV.19.6.3}
With the preceding notations, suppose that the suite $\mathbf{f}$ is $\operatorname{gr}_.(N)$-regular.
Consider the following properties:
\begin{enumerate}[label=\emph{\alph*)}]
  \item $\mathbf{f}$ is $\operatorname{gr}_.(M)$-regular.
  \item $\psi_M$ is bijective.
  \item $\psi'_R$ is surjective (in other words $\operatorname{gr}'_.(R)$ is generated as $\operatorname{gr}^*_{\mathfrak{J}}(A)$-module by $\operatorname{Im}(\operatorname{gr}_.(R)\to\operatorname{gr}'_.(R))$).
  \item $\psi'_R$ is injective.
  \item[d')] The homomorphism $(\psi'_R)_0:\operatorname{gr}_.(R)\otimes_A(A/\mathfrak{J})\to\operatorname{gr}'_.(R)$ obtained by restricting $\psi'_R$ to polynomials of degree $0$, is injective.
  \item $\psi'_R$ is bijective.
  \item[f)] $j'$ is injective.
  \item[f')] The homomorphism $\operatorname{gr}_.(R)\otimes_A(A/\mathfrak{J})\to\operatorname{gr}_.(N)\otimes_A(A/\mathfrak{J})$ obtained by restricting $j'$ to polynomials of degree $0$, is injective.
  \item The $\mathfrak{J}$-preadic filtration of $\operatorname{gr}_.(R)$ is induced by that of $\operatorname{gr}_.(N)$.
  \item The filtrations $(R'_k)$ and $(R''_k)$ on $R$ are identical.
\end{enumerate}
One has the implications:
\[
  a)\Rightarrow e)\Rightarrow b)\Leftrightarrow c)\Leftrightarrow h)
\]
\[
  g)\Rightarrow d)\Leftrightarrow d')\Leftrightarrow f)\Leftrightarrow f')
\]
Moreover, when every module quotient of a sub-module of $\operatorname{gr}_k(M)$ is separated for the $\mathfrak{J}$-preadic topology (which will be the case when $A$ is Noetherian, the $f_i$ belong to the radical of $A$ and the $\operatorname{gr}_k(M)$ are A-modules of finite type), \emph{then all the conditions \emph{a)} to \emph{h)} are equivalent}.
\end{proposition}

\begin{corollary}[19.6.4]
\label{IV.19.6.4}
\oldpage[IV]{211}
Let $\mathfrak{K}$ be an ideal of $A$, $\mathfrak{L}=\mathfrak{J}+\mathfrak{K}$.
Suppose that the filtration $(N_k)$ (hence also $(M_k)$ and $(N'_k)$) is the $\mathfrak{K}$-preadic filtration.
Consider $\operatorname{gr}_.(R)$ (resp.\ $\operatorname{gr}'_.(R)$) as a graded module over $\operatorname{gr}^*_{\mathfrak{K}}(A)$ (resp.\ $\operatorname{gr}^*_{\mathfrak{L}}(A)$).
\begin{enumerate}[label=\emph{(\roman*)}]
  \item Let $S$ be a part of $\operatorname{gr}_.(R)$ which generates $\operatorname{gr}_.(R)$ as $\operatorname{gr}^*_{\mathfrak{K}}(A)$-module.
  Then the condition $c)$ of (\sref{IV.19.6.3}) is equivalent to the following condition:
  \begin{enumerate}[label=\emph{c')}]
    \item The image of $S$ in $\operatorname{gr}'_.(R)$ by $\psi'_R$ generates $\operatorname{gr}'_.(R)$ as $\operatorname{gr}^*_{\mathfrak{L}}(A)$-module.
  \end{enumerate}
  \item Suppose verified the condition \emph{e)} of (\sref{IV.19.6.3}), and suppose moreover that the quotients of the $\operatorname{gr}_k(R)$ are separated for the $\mathfrak{J}$-preadic topology (which will be the case when $A$ is Noetherian, the $f_i$ are in the radical of $A$ and the $\operatorname{gr}_k(M)$ are A-modules of finite type).
  Then, for a part $S$ of homogeneous elements to generate $\operatorname{gr}_.(R)$ as $\operatorname{gr}^*_{\mathfrak{K}}(A)$-module, it is necessary and sufficient that its image by $\psi'_R$ generate $\operatorname{gr}'_.(R)$ as $\operatorname{gr}^*_{\mathfrak{L}}(A)$-module.
\end{enumerate}
\end{corollary}

\subsection{Hironaka's criterion of normal flatness.}
\label{subsection:IV.19.7}

\oldpage[IV]{212}

\begin{theorem}[19.7.1]
\label{IV.19.7.1}
(Hironaka.) ---
Let $A$ be a ring, $\mathfrak{K}$ an ideal of $A$, $\mathbf{f}=(f_1,\ldots,f_n)$ a suite of elements of $A$ which is $(A/\mathfrak{K})$-regular, $\mathfrak{J}=f_1 A+\ldots+f_n A$ the ideal generated by $\mathbf{f}$, $\mathfrak{L}=\mathfrak{J}+\mathfrak{K}$, $M$ an A-module.
One has in this case $(\operatorname{gr}^*_{\mathfrak{K}}(M))\otimes_A(A/\mathfrak{J})=(\operatorname{gr}^*_{\mathfrak{K}}(M))\otimes_A(A/\mathfrak{L})=(\operatorname{gr}^*_{\mathfrak{K}}(M))\otimes_{A/\mathfrak{K}}(A/\mathfrak{L})$, since $(\mathfrak{K}^n M/\mathfrak{K}^{n+1}M)\otimes_A(A/\mathfrak{J})=\mathfrak{K}^n M/(\mathfrak{K}^{n+1}+\mathfrak{J}\mathfrak{K}^n)M=\mathfrak{K}^n M/(\mathfrak{K}^{n+1}+\mathfrak{L}\mathfrak{K}^n)M$.
Consider the following conditions:
\begin{enumerate}[label=\emph{\alph*)}]
  \item $\operatorname{gr}^*_{\mathfrak{K}}(M)$ is a flat $(A/\mathfrak{K})$-module.
  \item $\operatorname{gr}^*_{\mathfrak{L}}(M)$ is a flat $(A/\mathfrak{L})$-module and the canonical homomorphism ($\mathbf{0}$, \sref{IV.15.1.5.2})
  \[
    \psi_M:((\operatorname{gr}^*_{\mathfrak{K}}(M))\otimes_A(A/\mathfrak{L}))[\mathrm{T}_1,\ldots,\mathrm{T}_n]\to\operatorname{gr}^*_{\mathfrak{L}}(M)
  \]
  is bijective.
  \item $(\operatorname{gr}^*_{\mathfrak{K}}(M))\otimes_{A/\mathfrak{K}}(A/\mathfrak{L})$ is a flat $(A/\mathfrak{L})$-module and the suite $\mathbf{f}$ is $(\operatorname{gr}^*_{\mathfrak{K}}(M))$-regular.
  \item $\operatorname{gr}^*_{\mathfrak{L}}(M)$ is a flat $(A/\mathfrak{L})$-module, and for every $\mathfrak{p}\in\operatorname{Ass}_{A/\mathfrak{K}}(A/\mathfrak{L})$, $(\operatorname{gr}^*_{\mathfrak{K}}(M))_\mathfrak{p}$ is a flat $(A/\mathfrak{K})_\mathfrak{p}$-module.
\end{enumerate}
One then has the implications
\[
  a)\Rightarrow c)\Rightarrow b)
\]
and $a)\Rightarrow d)$.

When every quotient of a sub-module of $\operatorname{gr}^p_{\mathfrak{K}}(M)$ ($p\geq 0$) is ideally separated for $\mathfrak{J}$ (Bourbaki, \emph{Alg.~comm.}, chap.~III, \S~5, n\textsuperscript{o}~1), the conditions \emph{a)}, \emph{b)} and \emph{c)} are equivalent.

Finally, when $A$ is Noetherian, the $f_i$ belong to the radical of $A$, the suite $\mathbf{f}$ is $\operatorname{gr}^*_{\mathfrak{K}}(A)$-regular and $M$ is an A-module of finite type, the conditions \emph{a)}, \emph{b)}, \emph{c)}, \emph{d)} are equivalent.
\end{theorem}

\begin{remark}[19.7.2]
\label{IV.19.7.2}
\oldpage[IV]{218}
\emph{(i)} When one supposes that $\operatorname{gr}^*_{\mathfrak{K}}(A)$ is a flat $(A/\mathfrak{K})$-module, the hypothesis that the suite $\mathbf{f}$ is $(A/\mathfrak{K})$-regular implies that it is also $\operatorname{gr}^*_{\mathfrak{K}}(A)$-regular ($\mathbf{0}$, \sref{IV.15.1.14}).
One notes in particular that this will be the case when $A$ and $A/\mathfrak{K}$ \emph{are regular rings} ($\mathbf{0}$, \sref{IV.17.3.6}): indeed, $\mathfrak{K}$ is then a regular ideal of $A$ (\sref{IV.19.1.2}), hence quasi-regular, and localising at the maximal ideals of $A$ containing $\mathfrak{K}$ and using ($\mathbf{0}$, \sref{IV.15.1.7}), one sees that $\operatorname{gr}^*_{\mathfrak{K}}(A)$ is a flat $(A/\mathfrak{K})$-module \emph{projective}.

\emph{(ii)} It would be interesting to be able to prove the last conclusion of (\sref{IV.19.7.1}) by only supposing that there exists a B-algebra $B$ which is a Noetherian ring, for which $\mathfrak{J}B$ is contained in the radical of $B$, and $M$ is a B-module of finite type.
\end{remark}

\begin{corollary}[19.7.3]
\label{IV.19.7.3}
Let $A$ be a local Noetherian ring, $\mathfrak{m}$ its maximal ideal, $\mathfrak{K}$ an ideal of $A$ such that the ring $A/\mathfrak{K}$ is regular and of dimension $n$, $M$ an A-module of finite type.
The following conditions are equivalent:
\begin{enumerate}[label=\emph{\alph*)}]
  \item $\operatorname{gr}^*_{\mathfrak{K}}(M)$ is a flat $(A/\mathfrak{K})$-module.
  \item If $P(T)$ and $Q(T)$ are the Poincar\'e series of the $(A/\mathfrak{m})$-graded modules $\operatorname{gr}^*_{\mathfrak{m}}(M)$ and $\operatorname{gr}^*_{\mathfrak{K}}(M)\otimes_{A/\mathfrak{K}}(A/\mathfrak{m})$, one has
  \begin{equation}
  \label{IV.19.7.3.1}
    Q(T)=P(T)(1-T)^{-n}.
  \end{equation}
\end{enumerate}
\end{corollary}

\begin{corollary}[19.7.4]
\label{IV.19.7.4}
Let $X$ be a locally Noetherian prescheme, $Y$ a closed sub-prescheme of $X$, regular and connected, $\mathscr{F}$ an $\sh{O}_X$-Module coherent normally flat along $Y$ (\sref{IV.6.10.1}).
Then there exists a formal power series $R\in\mathbb{Q}[[T]]$ (independent of $x\in Y$) such that, for every $x\in Y$, the Poincar\'e series of $\operatorname{gr}^*_{\mathfrak{m}_x}(\mathscr{F}_x)$ is given by the formula
\begin{equation}
\label{IV.19.7.4.1}
  P_x(\mathscr{F})(T)=R(T)(1-T)^{-n}\quad\text{where }n=\dim(\sh{O}_{X,x}).
\end{equation}
\end{corollary}

\begin{corollary}[19.7.5]
\label{IV.19.7.5}
\oldpage[IV]{219}
Let $X$ be a locally Noetherian prescheme, $Y$ a closed sub-prescheme of $X$, $Z$ a closed sub-prescheme of $Y$; suppose $Y$ and $Z$ regular.
Let $\mathscr{F}$ be an $\sh{O}_X$-Module coherent.
\begin{enumerate}[label=\emph{(\roman*)}]
  \item If $\mathscr{F}$ is normally flat along $Y$ at the points of $Z$ (cf.\ (\sref{IV.11.3.4}) for the definition of normal flatness at a point), \emph{then $\mathscr{F}$ is normally flat along the points of $Z$}.
  \item Suppose moreover that $Z$ is connected and that $\sh{O}_X$ is normally flat along $Y$ at the points of $Z$ (which will take place in particular if $X$ is regular at the points of $Z$).
  \emph{If $\mathscr{F}$ is normally flat along $Z$}, and if there exists a point $z\in Z$ such that $\mathscr{F}$ is normally flat along $Y$ at the point $z$, \emph{then $\mathscr{F}$ is normally flat along $Y$ at all points of $Z$}.
\end{enumerate}
\end{corollary}

\subsection{Properties of passage to the projective limit.}
\label{subsection:IV.19.8}

\oldpage[IV]{219}
In this section, the notations and conventions on projective limits are those of (\sref{IV.8.5.1}) and (\sref{IV.8.1.1}).

\begin{proposition}[19.8.1]
\label{IV.19.8.1}
Suppose that the transition morphisms $S_\mu\to S_\lambda$ ($\lambda\leq\mu$) are flat, and in addition that one of the two following hypotheses is verified:
\begin{enumerate}
  \item[1\textsuperscript{o}] The preschemes $S_\lambda$ are locally Noetherian.
  \item[2\textsuperscript{o}] The transition morphisms $S_\mu\to S_\lambda$ are surjective (hence faithfully flat).
\end{enumerate}

In these conditions:
\begin{enumerate}[label=\emph{(\roman*)}]
  \item Suppose $S_\alpha$ quasi-compact.
  Let $\mathscr{F}_\alpha$ be an $\sh{O}_{S_\alpha}$-Module quasi-coherent, which we suppose moreover of finite type when hypothesis $1\textsuperscript{o}$ is satisfied.
  Let $(f_{i\alpha})_{1\leq i\leq n}$ be a finite suite of sections of $\mathscr{F}_\alpha$ above $S_\alpha$.
  For the suite of sections $f_i$ of the $\sh{O}_S$-Module $\mathscr{F}$ above $S$, corresponding to the $f_{i\alpha}$, to be $\mathscr{F}$-regular, it is necessary and sufficient that there exist $\lambda\geq\alpha$ such that the suite of sections $f_{i\lambda}$ of $\mathscr{F}_\lambda$ above $S_\lambda$ be $\mathscr{F}_\lambda$-regular.
  \item Suppose $X_\alpha$ quasi-compact, and let $j_\alpha:Y_\alpha\to X_\alpha$ a locally finite immersion, which we suppose moreover locally of finite presentation when hypothesis $2\textsuperscript{o}$ is satisfied.
  Then, for the corresponding immersion $j:X\to S$ to be regular, it is necessary and sufficient that there exist $\lambda\geq\alpha$ such that $j_\lambda:X_\lambda\to S_\lambda$ be regular.
\end{enumerate}
\end{proposition}

\begin{proposition}[19.8.2]
\label{IV.19.8.2}
\oldpage[IV]{220}
Let $X_\alpha$ be an $S_\alpha$-prescheme quasi-compact and locally of finite presentation.
\begin{enumerate}[label=\emph{(\roman*)}]
  \item Let $\mathscr{F}_\alpha$ be an $\sh{O}_{X_\alpha}$-Module quasi-coherent of finite presentation, $(f_{i\alpha})_{1\leq i\leq n}$ a suite of sections of $\sh{O}_{X_\alpha}$ above $X_\alpha$.
  For the suite of sections $f_i$ of $\sh{O}_X$ above $X$, corresponding to the $f_{i\alpha}$, to be $\mathscr{F}_X$-transversally regular relatively to $S$ (\sref{IV.19.2.1}), it is necessary and sufficient that there exist $\lambda\geq\alpha$ such that the suite of sections $f_{i\lambda}$ of $\sh{O}_{X_\lambda}$ above $X_\lambda$ be $\mathscr{F}_\lambda$-transversally regular relatively to $S_\lambda$.
  \item Let $Y_\alpha$ be an $S_\alpha$-prescheme quasi-compact, $j_\alpha:Y_\alpha\to X_\alpha$ an $S_\alpha$-immersion locally of finite presentation.
  For the S-immersion corresponding $j:Y\to X$ to be transversally regular relatively to $S$, it is necessary and sufficient that there exist $\lambda\geq\alpha$ such that $j_\lambda:Y_\lambda\to X_\lambda$ be transversally regular relatively to $S_\lambda$.
\end{enumerate}
\end{proposition}

\begin{remark}[19.8.3]
\label{IV.19.8.3}
The preceding proofs show that the statements of (\sref{IV.19.8.1}) and (\sref{IV.19.8.2}) subsist when one replaces the conditions concerning regularity or transverse regularity in $X$ (resp.\ $X_\lambda$) by these conditions in a neighbourhood of a point $x\in X$ (resp.\ in a neighbourhood of $x_\lambda$, projection of $x$).
\end{remark}

